%=============================================================================
% KAPPA OBSERVATIONAL BOUNDS
% What do existing observations constrain the constitutive split parameter to?
% Can experiment resolve the two-value ambiguity (0.27 vs 0.54)?
%
% Date: March 23, 2026
% Status: CRITICAL FINDING -- kappa is essentially UNCONSTRAINED by current data
%=============================================================================
\documentclass[11pt,a4paper]{article}
\usepackage[margin=1in]{geometry}
\usepackage{amsmath,amssymb,amsfonts,amsthm}
\usepackage{booktabs}
\usepackage{hyperref}
\usepackage{xcolor}
\usepackage{tcolorbox}

\newtheorem{theorem}{Theorem}
\newtheorem{proposition}{Proposition}
\newtheorem{finding}{Finding}
\theoremstyle{definition}
\newtheorem{definition}{Definition}

\title{Observational Bounds on the Constitutive Split Parameter~$\kappa$:\\
Can Experiment Resolve the Two-Value Ambiguity?}

\author{DFD Research Campaign --- March 23, 2026}
\date{}

\begin{document}
\maketitle

%=============================================================================
\section{Executive Summary}
%=============================================================================

\begin{tcolorbox}[colback=red!5!white, colframe=red!60!black,
  title=Principal Finding]
\textbf{The constitutive split parameter $\kappa$ is essentially
unconstrained by all existing observational data.}
Both theoretical values ($\kappa = 3/11 \approx 0.27$ and
$\kappa = \cos 2\theta_W \approx 0.54$) are equally consistent with
every published measurement.  The reason is structural:
DFD's constitutive split preserves $v_{\rm ph} = c/n$ for
\emph{both} polarizations, so the stringent astrophysical
birefringence bounds that might otherwise kill one value
\emph{do not apply at all}.
\end{tcolorbox}

\begin{tcolorbox}[colback=green!5!white, colframe=green!60!black,
  title=Resolution Path]
Only a \textbf{dedicated TE/TM cavity experiment}---comparing
pure-TE and pure-TM modes in an SRF resonator while modulating the
gravitational potential (annual solar modulation or altitude
change)---can distinguish the two values.  The predicted splitting is:
\begin{equation}
\frac{\delta\nu_{\rm TE} - \delta\nu_{\rm TM}}{\nu}
= 2\kappa\,\Delta\psi,
\end{equation}
which for the annual solar modulation
($\Delta\psi_\odot \approx 3.3 \times 10^{-10}$) gives
$3.6 \times 10^{-10}$ (EW route) or $1.8 \times 10^{-10}$ (EH route)
---a factor-of-two difference that current SRF technology can resolve.
\end{tcolorbox}


%=============================================================================
\section{The Two-Value Problem}
\label{sec:two-values}
%=============================================================================

The DFD constitutive split parameter $\kappa$ governs the
asymmetric response of vacuum permittivity and permeability to
the density field $\psi$:
\begin{equation}\label{eq:constitutive-split}
\epsilon(\psi) = \epsilon_0\, n\, e^{+\kappa\psi}, \qquad
\mu(\psi) = \mu_0\, n\, e^{-\kappa\psi},
\end{equation}
where $n = e^\psi$.  The product $\epsilon\mu = \epsilon_0\mu_0 n^2$
is preserved, guaranteeing that the phase velocity
$v_{\rm ph} = 1/\sqrt{\epsilon\mu} = c/n$ is
\emph{polarization-independent}.

Two independent derivation routes yield different values:

\subsection{Route 1: Euler--Heisenberg (QED sector only)}

The Euler--Heisenberg effective Lagrangian for photon--photon
scattering in the weak-field limit is:
\begin{equation}
\mathcal{L}_{\rm EH} = \frac{\alpha}{90\pi B_c^2}
\left[4(\mathbf{E}^2 - \mathbf{B}^2)^2
+ 7(\mathbf{E}\cdot\mathbf{B})^2\right].
\end{equation}
The asymmetric coefficient ratio $4:7$ between the $\mathbf{E}$
and $\mathbf{B}$ sectors gives:
\begin{equation}
\boxed{\kappa_{\rm EH} = \frac{7 - 4}{7 + 4} = \frac{3}{11}
\approx 0.273.}
\end{equation}
This route probes only the QED vacuum polarization.

\subsection{Route 2: Electroweak mixing (full EW sector)}

The photon inherits its coupling to $\psi$ from the electroweak
$B$ and $W^3$ bosons, with hypercharge coupling $e\sin\theta_W$
and isospin coupling $e\cos\theta_W$ respectively.  After
symmetry breaking:
\begin{equation}
\boxed{\kappa_{\rm EW} = \frac{\cos^2\theta_W - \sin^2\theta_W}
{\cos^2\theta_W + \sin^2\theta_W} = \cos 2\theta_W
= 1 - 2\sin^2\theta_W \approx 0.538.}
\end{equation}
With $\sin^2\theta_W = 0.2312$ (PDG 2024).

\subsection{Route 3: Trace anomaly (sign only)}

The QED trace anomaly $\langle T^\mu{}_\mu\rangle =
(\alpha/6\pi)(\mathbf{B}^2 - \mathbf{E}^2)$ shows the
$\mathbf{B}$ sector dominates, implying $\kappa > 0$, but does
not determine the magnitude.

\subsection{Physical distinction}

The two routes differ because they probe different levels of the
vacuum structure:
\begin{itemize}
\item $\kappa_{\rm EH}$: the $\psi$ field couples to the
  \emph{QED vacuum} only (electron loops)
\item $\kappa_{\rm EW}$: the $\psi$ field couples to the
  \emph{full electroweak vacuum} (W/Z/Higgs)
\end{itemize}
DFD favours the electroweak route because $\psi$ is a
gravitational-strength scalar that couples universally.


%=============================================================================
\section{Why Most Observational Bounds Do Not Apply}
\label{sec:inapplicable}
%=============================================================================

\begin{tcolorbox}[colback=yellow!5!white, colframe=yellow!60!black,
  title=Critical Structural Point]
The DFD constitutive split preserves
$v_{\rm ph} = c/n$ for \textbf{both} polarizations
(Eq.~\ref{eq:constitutive-split}).  This is \emph{not} vacuum
birefringence in the standard sense.  The split affects
impedance $Z = \sqrt{\mu/\epsilon}$, not phase velocity.
\end{tcolorbox}

\subsection{Astrophysical birefringence bounds: do NOT apply}

The most stringent tests of vacuum birefringence come from:
\begin{itemize}
\item \textbf{GRB polarization:} Gamma-ray burst afterglows
  constrain Lorentz-violating birefringence coefficients at
  $|\tilde\kappa_{\rm LV}| < 10^{-34}$ (SME framework).
\item \textbf{CMB polarization:} Cosmic birefringence
  (``cosmic polarization rotation'') bounded at $< 0.1^\circ$.
\item \textbf{AGN jet polarization:} VLBI observations
  constrain polarization rotation at milliarcsecond scales.
\end{itemize}
\textbf{None of these apply to DFD's $\kappa$.}  All test
whether different polarizations propagate at different
\emph{speeds}, i.e., whether $v_{\rm ph}^{\rm TE} \neq
v_{\rm ph}^{\rm TM}$.  In DFD, $v_{\rm ph}^{\rm TE} =
v_{\rm ph}^{\rm TM} = c/n$ exactly, by construction.

\subsection{Shapiro delay polarization: does NOT apply}

The Cassini measurement of the PPN parameter $\gamma$
(Bertotti, Iess \& Tortora 2003) achieved
$\gamma - 1 = (2.1 \pm 2.3) \times 10^{-5}$ from the
Shapiro delay of radio signals passing near the Sun.
The measurement used dual-frequency (X-band and Ka-band)
tracking, but was not sensitive to polarization-dependent
delay because:
\begin{enumerate}
\item DFD predicts $v_{\rm ph}$ is polarization-independent.
\item The Shapiro delay $\Delta t = (2GM/c^3)\ln(4r_1 r_2/b^2)
  \times (1 + \gamma)/2$ depends on $v_{\rm ph}$, not on $Z$.
\item DFD's PPN $\gamma = 1$ exactly (the optical metric is
  conformally flat), consistent with Cassini.
\end{enumerate}

\subsection{VLBI solar deflection: does NOT apply}

VLBI measurements of radio source deflection by the Sun
($\gamma = 0.9998 \pm 0.0004$) constrain the conformal
coupling $\psi$, not the constitutive split.  Photon
trajectories follow null geodesics of the \emph{same} optical
metric regardless of polarization.

\subsection{SME photon-sector bounds: do NOT apply}

The Standard-Model Extension (SME) photon sector parameterizes
Lorentz-violating terms via $\tilde\kappa$ coefficients.  The
birefringent sector has $|\tilde\kappa| < 10^{-34}$ from GRB
observations.  However, the DFD constitutive split:
\begin{enumerate}
\item Does \emph{not} violate Lorentz invariance---it couples
  to a scalar field $\psi$ that respects local Lorentz symmetry.
\item Does \emph{not} produce birefringence in the standard sense
  (different phase velocities).
\item Maps onto the \emph{birefringence-free} sector of the
  SME, which has much weaker bounds
  ($|\hat\kappa_{\rm tr}| \lesssim 10^{-15}$ to $10^{-20}$
  for dimension-4 coefficients).
\end{enumerate}
Even the birefringence-free SME bounds do not directly constrain
DFD's $\kappa$, because the DFD effect is modulated by the
gravitational potential $\psi$ rather than being a constant
background coefficient.

\subsection{Drummond--Hathrell QED gravitational birefringence:
  not relevant}

The Drummond--Hathrell (1980) prediction of QED gravitational
birefringence arises from one-loop vacuum polarization on curved
spacetime.  Its magnitude near the Sun is of order:
\begin{equation}
\Delta n_{\rm DH} \sim \alpha\,
\left(\frac{\bar\lambda_C}{R_\odot}\right)^{\!2}
\frac{r_s}{R_\odot}
\sim 10^{-50},
\end{equation}
where $\bar\lambda_C = \hbar/(m_e c) \approx 3.86 \times 10^{-13}$~m
is the electron Compton wavelength and
$r_s = 2GM_\odot/c^2 \approx 2950$~m.
This is a different physical effect (tidal coupling to curvature)
and is unmeasurably small.


%=============================================================================
\section{Bounds That Do Apply (and Why They Are Weak)}
\label{sec:applicable}
%=============================================================================

\subsection{Bound 1: Cavity parametric stability}
\label{sec:cavity-stability}

Appendix R of the DFD Unified Theory establishes that
$|\kappa| \lesssim 1$ from the absence of $2\omega$ parametric
instabilities in extreme-$Q$ resonators.  If $\kappa$ were much
larger than unity, the coupling $(\kappa/2)\dot\psi\,\mathcal{B}$
in the Poynting theorem would pump energy between the EM and
$\psi$ sectors at observable rates.

\textbf{Constraint:} $|\kappa| \lesssim 1$.  Both 0.27 and 0.54
trivially satisfy this.

\subsection{Bound 2: Clock comparisons (LPI tests)}
\label{sec:clock}

The best clock comparison data come from:
\begin{center}
\small
\begin{tabular}{lcl}
\toprule
\textbf{Experiment} & \textbf{Bound} & \textbf{Reference} \\
\midrule
H-maser/Cs (14-yr) & $|\beta_{\rm H} - \beta_{\rm Cs}| < 2.5
  \times 10^{-7}$ & Ashby et al.\ (2018) \\
Yb$^+$ E3/E2 (PTB) & $|\beta_{\rm E3} - \beta_{\rm E2}| <
  10^{-8}$ & Lange et al.\ (2021) \\
Al$^+$/Hg$^+$ & $|\beta_{\rm Al} - \beta_{\rm Hg}| < 5.3
  \times 10^{-7}$ & Rosenband et al.\ (2008) \\
BACON (Al$^+$/Sr/Yb) & Multi-species at $8 \times 10^{-18}$ &
  Beloy et al.\ (2021) \\
\bottomrule
\end{tabular}
\end{center}

In DFD, the relevant coupling is (Appendix~R):
\begin{equation}
\xi_{\rm LPI}^{\rm res}(\kappa)
= \text{(screened residual of)}\;
  1 - \alpha_L^{(M)} - K_\epsilon^{(S)}\kappa
  + \mathcal{O}(\kappa^2).
\end{equation}
After the constitutive-chain cancellation (Sec.~10 of the Unified
Theory), the leading geometric effect vanishes at tree level.
The surviving $\kappa$-dependent signal is:
\begin{equation}
\Delta\xi(\kappa) \sim \kappa \times
\underbrace{K_\epsilon^{(S)}}_{\sim 1\text{--}3}
\times \underbrace{(\text{loop correction})}_{\sim \alpha/(2\pi)}
\times \underbrace{\Delta\psi_\odot}_{\sim 3.3\times 10^{-10}}.
\end{equation}
For $\kappa = 0.54$:
$\Delta\xi \sim 0.54 \times 2 \times (1/860) \times 3.3
\times 10^{-10} \sim 4 \times 10^{-13}$.

Even the tightest LPI bound ($10^{-8}$ from PTB Yb$^+$)
gives:
\begin{equation}
\kappa < \frac{10^{-8}}{4 \times 10^{-13}/\kappa}
\sim 2.5 \times 10^4.
\end{equation}
\textbf{This is not constraining at all.}  The tree-level
cancellation shields $\kappa$ by many orders of magnitude.

\subsection{Bound 3: Cassini PPN $\gamma$}
\label{sec:cassini}

Cassini constrains $|\gamma - 1| < 2.3 \times 10^{-5}$.
In DFD, $\gamma = 1$ exactly from the conformal structure of
the optical metric, independent of $\kappa$.  The constitutive
split does not modify photon trajectories or the Shapiro delay.

\textbf{Constraint on $\kappa$:} None.

\subsection{Bound 4: Gravitational wave polarization}
\label{sec:gw}

LIGO/Virgo measurements constrain the speed difference between
gravitational wave polarizations at
$|v_+ - v_\times|/c < 10^{-15}$.  This constrains modified
gravity theories but \emph{not} the electromagnetic constitutive
split, which is a photon-sector parameter.

\textbf{Constraint on $\kappa$:} None.


%=============================================================================
\section{Quantitative Predictions for Each $\kappa$ Value}
\label{sec:predictions}
%=============================================================================

\subsection{TE/TM cavity splitting}

For a cavity in a gravitational potential with
$\psi = GM/(Rc^2)$:

\begin{center}
\begin{tabular}{lccc}
\toprule
\textbf{Observable} & $\kappa = 0.54$ & $\kappa = 0.27$ &
  $\kappa = 0$ \\
\midrule
\multicolumn{4}{l}{\textit{At Sun's surface}
  ($\psi_\odot = 2.12 \times 10^{-6}$):} \\
$(\delta\nu_{\rm TE} - \delta\nu_{\rm TM})/\nu$ &
  $2.3 \times 10^{-6}$ & $1.2 \times 10^{-6}$ & $0$ \\[4pt]
\multicolumn{4}{l}{\textit{Annual solar modulation at Earth}
  ($\Delta\psi = 3.3 \times 10^{-10}$):} \\
$\Delta(\delta\nu_{\rm TE} - \delta\nu_{\rm TM})/\nu$ &
  $3.6 \times 10^{-10}$ & $1.8 \times 10^{-10}$ & $0$ \\[4pt]
\multicolumn{4}{l}{\textit{At Earth's surface (Earth potential)}
  ($\psi_\oplus = 6.9 \times 10^{-10}$):} \\
$(\delta\nu_{\rm TE} - \delta\nu_{\rm TM})/\nu$ &
  $7.5 \times 10^{-10}$ & $3.8 \times 10^{-10}$ & $0$ \\
\bottomrule
\end{tabular}
\end{center}

\textbf{Key discrimination:} The annual modulation signal differs by
a factor of two between the EW and EH routes.  SRF cavities with
$Q > 10^{11}$ and frequency stability $\sigma_y \sim 10^{-18}$ at
$10^4$~s integration can reach $10^{-10}$ fractional resolution,
making this measurement feasible.

\subsection{Impedance modulation}

The vacuum impedance acquires a $\psi$-dependent correction:
\begin{equation}
Z(\psi) = Z_0\,e^{-\kappa\psi},
\end{equation}
where $Z_0 = \sqrt{\mu_0/\epsilon_0} = 376.73$~$\Omega$.
For $\Delta\psi_\odot = 3.3 \times 10^{-10}$:
\begin{equation}
\frac{\Delta Z}{Z_0} = -\kappa\,\Delta\psi
= \begin{cases}
-1.8 \times 10^{-10} & (\kappa = 0.54), \\
-0.9 \times 10^{-10} & (\kappa = 0.27), \\
0 & (\kappa = 0).
\end{cases}
\end{equation}
This is far below current reflectometry precision
($\delta\Gamma/\Gamma \sim 10^{-6}$) but could be accessible
with SRF techniques ($\delta Z/Z \sim 10^{-12}$).

\subsection{Propagation birefringence}

\begin{equation}
\frac{v_{\rm ph}^{\rm TE} - v_{\rm ph}^{\rm TM}}{c} = 0
\qquad \text{for all }\kappa.
\end{equation}
\textbf{No propagation birefringence in DFD}, regardless of
$\kappa$.  This is a structural prediction, not a fine-tuning.

\subsection{LPI violation parameter}

\begin{equation}
\xi_{\rm LPI}^{\rm res}(\kappa) = K_\epsilon^{(S)}\kappa
\times (\text{screening factor}).
\end{equation}
At Earth's surface in the screened regime, the screening
factor is $\ll 1$ (constrained by BACON to give
$\xi_{\rm LPI}^{\rm res} < 10^{-5}$).  The
$\kappa$-dependence is buried under the screening.


%=============================================================================
\section{Is $\kappa = 0$ Experimentally Excluded?}
\label{sec:kappa-zero}
%=============================================================================

No.  There is currently \textbf{no observational evidence} that
$\kappa \neq 0$.  The parameter $\kappa = 0$ corresponds to the
symmetric Tamm--Plebanski constitutive relations
$\varepsilon_{\rm eff} = \varepsilon_0 e^{+\psi}$,
$\mu_{\rm eff} = \mu_0 e^{+\psi}$, in which both sectors respond
identically to $\psi$.

However, $\kappa = 0$ is \textbf{theoretically disfavoured} by
both derivation routes:
\begin{enumerate}
\item The Euler--Heisenberg Lagrangian has an inherent $E/B$
  asymmetry (coefficient ratio $4:7$), implying $\kappa \neq 0$
  at the QED level.
\item The electroweak mixing angle $\sin^2\theta_W \neq 1/2$
  implies $\cos 2\theta_W \neq 0$, so $\kappa \neq 0$ at
  the EW level.
\item The QED trace anomaly has a definite sign
  ($\mathbf{B}^2 - \mathbf{E}^2$), implying $\kappa > 0$.
\end{enumerate}
Nevertheless, without a direct experimental measurement, one
cannot exclude $\kappa = 0$ on observational grounds.  The
constitutive-chain cancellation protects the theory from
catastrophic disagreement with data for \emph{any} value of
$\kappa \in [0, 1]$.


%=============================================================================
\section{The Tightest Current Bound}
\label{sec:tightest}
%=============================================================================

\begin{tcolorbox}[colback=blue!5!white, colframe=blue!50!black,
  title=Tightest Observational Bound on $\kappa$]
\textbf{Result:} $|\kappa| \lesssim 1$ from cavity parametric
stability (Appendix~R).

All other bounds are orders of magnitude weaker when applied
to the DFD constitutive split, because the tree-level cancellation
(Section~10) screens $\kappa$ from clock comparisons, and the
preserved phase velocity screens it from all propagation-based
tests.

The two-value ambiguity ($0.27$ vs $0.54$) \textbf{cannot be
resolved by existing data}.
\end{tcolorbox}


%=============================================================================
\section{How to Resolve the Ambiguity: Experimental Proposals}
\label{sec:proposals}
%=============================================================================

\subsection{Proposal 1: TE/TM dual-mode SRF cavity}

\paragraph{Concept.}
Operate two SRF cavity modes---one predominantly TE, one
predominantly TM---locked to a common reference.  Monitor the
TE--TM frequency difference over one year.

\paragraph{Signal.}
\begin{equation}
\frac{\delta\nu_{\rm TE} - \delta\nu_{\rm TM}}{\nu}
= 2\kappa\,\Delta\psi_\odot\,\cos\!\left(\frac{2\pi t}{T_{\rm yr}}
+ \phi_0\right),
\end{equation}
with amplitude:
\begin{equation}
A = 2\kappa \times 3.3 \times 10^{-10}
= \begin{cases}
3.6 \times 10^{-10} & (\kappa = 0.54), \\
1.8 \times 10^{-10} & (\kappa = 0.27), \\
0 & (\kappa = 0).
\end{cases}
\end{equation}

\paragraph{Sensitivity requirement.}
SRF cavities at SQMS/Fermilab achieve $Q > 10^{11}$ at 1.3~GHz.
Allan deviation $\sigma_y \sim 10^{-18}$ at $10^4$~s is feasible
with cryogenic temperature stabilization.  Over a one-year
integration, the annual modulation can be extracted at SNR~$>10$
for $\kappa > 0.1$.

\paragraph{Discrimination.}
The two predictions differ by a factor of 2 in amplitude.
With 10\% amplitude measurement uncertainty, the routes are
distinguished at $> 5\sigma$.

\subsection{Proposal 2: Multi-species clock comparison with
  species-resolved $\kappa$ sensitivity}

Different atomic species have different $K_\epsilon^{(S)}$ values
(the sensitivity to $\epsilon$-changes).  By comparing three or
more species, one can extract $\kappa$ independently of the
overall screening factor.

The required precision is $\sigma_y < 10^{-19}$ for the
annual modulation, which is at the frontier of next-generation
optical clocks (BACON-2, JILA 3D lattice).

\subsection{Proposal 3: Space-based cavity at varying
  solar distance}

A mission carrying a dual-mode SRF cavity to different
heliocentric distances would see $\psi$ vary by orders of
magnitude more than the terrestrial annual signal.
At 0.1~AU: $\psi \sim 10^{-7}$, giving
$\Delta\nu_{\rm TE-TM}/\nu \sim 10^{-7}\kappa$---easily
measurable even with modest cavity $Q$.


%=============================================================================
\section{Implications for the DFD Programme}
\label{sec:implications}
%=============================================================================

\begin{enumerate}
\item \textbf{The two-value problem is \emph{not} a crisis.}
  Both values are observationally viable.  The ambiguity is
  a theoretical question (QED vs EW coupling level) that
  experiment can resolve cleanly.

\item \textbf{The TE/TM cavity experiment has triple value:}
  \begin{enumerate}
  \item Measures $\kappa$ directly;
  \item If $\kappa \neq 0$ is found, confirms the constitutive
    split and rules out $\kappa = 0$;
  \item The measured value identifies which vacuum sector
    (QED or EW) $\psi$ couples to, which has deep implications
    for the DFD quantum gravity programme.
  \end{enumerate}

\item \textbf{If $\kappa = \cos 2\theta_W$, then $\kappa$ runs.}
  At different energy scales, $\sin^2\theta_W$ varies
  ($\sin^2\theta_W(M_Z) = 0.2312$;
  $\sin^2\theta_W(0) \approx 0.2387$), so
  $\kappa(0) \approx 0.523$ vs $\kappa(M_Z) \approx 0.538$.
  This running is not detectable with current precision but
  is a distinctive prediction.

\item \textbf{$\kappa = 0$ is allowed but unlikely.}
  It would require both the EH asymmetry and the EW mixing
  to somehow not imprint on the constitutive relations---a
  fine-tuning that seems unnatural.
\end{enumerate}


%=============================================================================
\section{Summary Table: All Observational Channels}
\label{sec:summary}
%=============================================================================

\begin{table}[h!]
\centering
\caption{Complete inventory of observational channels and their
  constraining power on $\kappa$.}
\label{tab:summary}
\small
\begin{tabular}{p{4cm}ccp{4cm}}
\toprule
\textbf{Channel} & \textbf{Bound} &
  \textbf{Constrains $\kappa$?} & \textbf{Reason} \\
\midrule
GRB/AGN birefringence & $10^{-34}$ & \textbf{No} &
  Tests $v_{\rm ph}$ splitting; DFD has none \\
CMB polarization rotation & $<0.1^\circ$ & \textbf{No} &
  Same reason \\
Cassini Shapiro delay & $2.3\times10^{-5}$ & \textbf{No} &
  Constrains $\gamma$; $\kappa$ does not affect $v_{\rm ph}$ \\
VLBI deflection & $4.5\times10^{-4}$ & \textbf{No} &
  Constrains $\gamma$, not $Z$ \\
GW polarization speed & $10^{-15}$ & \textbf{No} &
  Wrong sector (GW, not EM) \\
\midrule
Cavity stability & $|\kappa| \lesssim 1$ & \textbf{Yes} &
  $2\omega$ instability absent \\
Clock comparisons & $\sim 10^{-8}$ & \textbf{Marginally} &
  Screened by tree-level cancellation \\
PTB Yb$^+$ E3/E2 & $10^{-8}$ & \textbf{Marginally} &
  Same-ion; probes residual only \\
\midrule
TE/TM cavity (proposed) & $10^{-10}$ & \textbf{YES} &
  \textit{Direct probe of $\kappa$} \\
Multi-species clocks & $10^{-19}$ & \textbf{Yes} &
  Requires next-gen clocks \\
Space cavity & $10^{-7}\kappa$ & \textbf{YES} &
  Large $\Delta\psi$ amplifies signal \\
\bottomrule
\end{tabular}
\end{table}


%=============================================================================
\section{Conclusions}
\label{sec:conclusions}
%=============================================================================

\begin{enumerate}
\item The constitutive split parameter $\kappa$ is
  \textbf{unconstrained} by all existing data at the
  level needed to distinguish between 0.27 and 0.54.

\item The structural reason is that DFD preserves
  $v_{\rm ph} = c/n$ for both polarizations, rendering
  the entire edifice of astrophysical birefringence
  bounds inapplicable.

\item The \textbf{only} applicable bound is
  $|\kappa| \lesssim 1$ from cavity parametric stability,
  which excludes neither value.

\item A \textbf{dedicated TE/TM SRF cavity experiment}
  can resolve the ambiguity by measuring the annual
  TE--TM frequency modulation at $\sim 10^{-10}$
  precision, where the two predictions differ by a
  factor of two.

\item $\kappa = 0$ is observationally allowed but
  theoretically disfavoured by both the Euler--Heisenberg
  asymmetry and the electroweak mixing angle.
\end{enumerate}


%=============================================================================
% REFERENCES
%=============================================================================

\begin{thebibliography}{20}

\bibitem{Bertotti2003}
B.~Bertotti, L.~Iess, and P.~Tortora,
``A test of general relativity using radio links with the Cassini spacecraft,''
Nature \textbf{425}, 374 (2003).

\bibitem{Lange2021}
R.~Lange et al.,
``Improved limits for violations of local position invariance from atomic clock comparisons,''
Phys.\ Rev.\ Lett.\ \textbf{126}, 011102 (2021).

\bibitem{Beloy2021}
K.~Beloy et al.\ (BACON Collaboration),
``Frequency ratio measurements at 18-digit accuracy using an optical clock network,''
Nature \textbf{591}, 564 (2021).

\bibitem{Ashby2018}
N.~Ashby et al.,
``A test of the gravitational redshift with Galileo satellites in an eccentric orbit,''
Phys.\ Rev.\ Lett.\ \textbf{121}, 231101 (2018).

\bibitem{Rosenband2008}
T.~Rosenband et al.,
``Frequency ratio of Al$^+$ and Hg$^+$ single-ion optical clocks; metrology at the 17th decimal place,''
Science \textbf{319}, 1808 (2008).

\bibitem{Drummond1980}
I.~T.~Drummond and S.~J.~Hathrell,
``QED vacuum polarization in a background gravitational field and its effect on the velocity of photons,''
Phys.\ Rev.\ D \textbf{22}, 343 (1980).

\bibitem{Euler1936}
H.~Euler and W.~Heisenberg,
``Folgerungen aus der Diracschen Theorie des Positrons,''
Z.~Phys.\ \textbf{98}, 714 (1936).

\bibitem{Kostelecky2011}
V.~A.~Kosteleck\'y and N.~Russell,
``Data tables for Lorentz and CPT violation,''
Rev.\ Mod.\ Phys.\ \textbf{83}, 11 (2011);
updated at \texttt{arxiv.org/abs/0801.0287}.

\bibitem{PDG2024}
R.~L.~Workman et al.\ (Particle Data Group),
Prog.\ Theor.\ Exp.\ Phys.\ \textbf{2024}, 083C01 (2024).

\bibitem{GravBiref2024}
Y.~Kubota et al.,
``Gravity-induced birefringence in spherically symmetric spacetimes,''
Phys.\ Rev.\ D \textbf{111}, 044001 (2025);
arXiv:2408.02729.

\bibitem{SMEbirefFree2026}
J.~J.~Wei et al.,
``Constraints on birefringence-free photon theory within standard-model extension,''
Phys.\ Rev.\ D (2026); arXiv:2602.01243.

\end{thebibliography}

\end{document}
